% Tercera lectura: conexiones con el resto del programa
\section{Coordenadas curvilíneas y bases locales}\label{cap:13}\label{sec:curvilineas}

Una parametrización
\[
  r:U\subset\R^2\longrightarrow\R^3,
  \qquad
  (u,v)\longmapsto r(u,v)
\]
describe una superficie mediante dos coordenadas. Los vectores
\[
  r_u=\frac{\partial r}{\partial u},
  \qquad
  r_v=\frac{\partial r}{\partial v}
\]
forman una base tangente local cuando son linealmente independientes.

\begin{center}
\begin{tikzpicture}[scale=.9]
  \draw[teallight,fill=teallight!60] (1,1)
    .. controls (3,.2) and (6,.5) .. (8,1.4)
    .. controls (7,3.8) and (3.2,4.5) .. (1,3)
    .. controls (.7,2.2) and (.7,1.6) .. cycle;
  \coordinate (p) at (4.6,2.1);
  \fill[ink] (p) circle (2.5pt) node[below left] {$p$};
  \draw[mainline,-{Latex[length=2.6mm]}] (p)--++(2.2,.25)
    node[right] {$r_u$};
  \draw[rust,line width=1.05pt,-{Latex[length=2.6mm]}] (p)--++(-.6,1.7)
    node[above] {$r_v$};
  \draw[muted,line width=.8pt,-{Latex[length=2.6mm]}] (p)--++(.2,2.2)
    node[above] {$n$};
  \node[text=muted] at (6.8,.8) {superficie};
\end{tikzpicture}
\end{center}

En coordenadas cartesianas la base es fija. En una superficie curva, la base cambia de
orientación y de longitud al cambiar el punto. Esta dependencia espacial es la razón por la
cual los productos internos y las derivadas adquieren coeficientes adicionales.

\section{Variedades y espacio tangente}\label{sec:variedades}

Una variedad es un objeto que localmente se parece a un espacio euclídeo. Una superficie
regular \(S\subset\R^3\) tiene en cada punto \(p\) un espacio tangente
\[
  T_pS=\operatorname{span}\{r_u(p),r_v(p)\}.
\]
El plano tangente es la mejor aproximación lineal local a la superficie. No es una proyección
cartográfica: pertenece al punto \(p\) y cambia al desplazarse sobre \(S\).

Si \(F:S\to M\) es una aplicación entre variedades, su derivada
\[
  dF_p:T_pS\longrightarrow T_{F(p)}M
\]
es una transformación lineal entre espacios tangentes. Esa transformación contiene la
información local sobre escala, orientación, deformación angular y deformación de áreas.

\section{Métrica y primera forma fundamental}\label{sec:metrica}

{\small\vspace{1mm}\linkUnidad{10_jacobiano_primera_forma}{Simulación interactiva: Unidad 14 — Jacobiano y primera forma fundamental}\par}


La métrica inducida por la parametrización queda determinada por
\[
  E=\langle r_u,r_u\rangle,\qquad
  F=\langle r_u,r_v\rangle,\qquad
  G=\langle r_v,r_v\rangle.
\]
El elemento de longitud es
\[
  \boxed{ds^2=E\,du^2+2F\,du\,dv+G\,dv^2}.
\]
En forma matricial:
\[
  ds^2=
  \begin{pmatrix}du&dv\end{pmatrix}
  \begin{pmatrix}E&F\\F&G\end{pmatrix}
  \begin{pmatrix}du\\dv\end{pmatrix}.
\]

\paragraph{Esfera.}
Para una esfera de radio \(R\), usando longitud \(\lambda\) y latitud \(\varphi\),
\[
  ds^2=R^2\cos^2\varphi\,d\lambda^2+R^2\,d\varphi^2.
\]
El factor \(\cos\varphi\) explica por qué un incremento fijo de longitud representa una
distancia cada vez menor al acercarse a los polos.

\section{Esfera, elipsoide y curvatura}

El elipsoide de revolución se caracteriza habitualmente mediante el semieje mayor \(a\) y
el achatamiento \(f\). La primera excentricidad satisface
\[
  e^2=2f-f^2.
\]
Dos radios de curvatura fundamentales son:
\[
  \boxed{
  M(\varphi)=
  \frac{a(1-e^2)}{(1-e^2\sin^2\varphi)^{3/2}}
  },
  \qquad
  \boxed{
  N(\varphi)=
  \frac{a}{\sqrt{1-e^2\sin^2\varphi}}
  }.
\]
\(M\) corresponde a la sección meridiana y \(N\) al primer vertical. La curvatura gaussiana
del elipsoide de revolución puede expresarse como
\[
  K=\frac{1}{MN}.
\]

\begin{center}
\begin{tikzpicture}[scale=.9]
  \draw[teal,line width=1.1pt] (5,2.5) ellipse (3.8 and 2.2);
  \draw[auxline] (1.2,2.5)--(8.8,2.5);
  \draw[auxline] (5,.3)--(5,4.7);
  \coordinate (p) at (7.7,3.75);
  \fill[rust] (p) circle (2.6pt) node[above right] {$p(\varphi)$};
  \draw[rust,line width=.8pt] ($(p)+(-1.7,.65)$)--($(p)+(1.7,-.65)$)
    node[right,font=\small] {tangente};
  \draw[mainline,-{Latex[length=2.5mm]}] (p)--++(-1.15,1.75)
    node[above] {normal};
  \node[text=muted] at (5,0) {elipsoide de revolución};
\end{tikzpicture}
\end{center}

\section{Proyecciones cartográficas}\label{sec:proyecciones}

Una proyección es una aplicación desde una superficie de referencia al plano:
\[
  P:S\longrightarrow\R^2.
\]
No puede conservar simultáneamente todas las propiedades métricas. Según la propiedad
priorizada se distingue:

\begin{center}
\begin{tabularx}{.92\textwidth}{>{\bfseries}l X X}
\toprule
Tipo & Conserva localmente & Ejemplos\\
\midrule
Conforme & ángulos & Mercator, Gauss--Krüger, UTM\\
Equivalente & áreas & Albers, Mollweide\\
Equidistante & ciertas distancias & azimutal equidistante\\
\bottomrule
\end{tabularx}
\end{center}

La derivada \(dP_p\) actúa entre planos tangentes. Sus valores singulares describen las
escalas principales de deformación; su determinante describe el cambio local de área.

\section{Variable compleja y conformidad}

Una función compleja \(w=f(z)\) es derivable en sentido complejo cuando existe
\[
  f'(z)=\lim_{\Delta z\to0}
  \frac{f(z+\Delta z)-f(z)}{\Delta z}.
\]
Cuando \(f'(z)\neq0\), la transformación lineal local es multiplicación por el complejo
\(f'(z)\): una rotación combinada con un escalado uniforme. Por ello las funciones holomorfas
preservan ángulos localmente.

La proyección estereográfica es un ejemplo central y conduce naturalmente al plano complejo
extendido
\[
  \widehat{\mathbb C}=\mathbb C\cup\{\infty\}.
\]
La relación entre rotaciones, multiplicación compleja y conformidad conecta la primera clase
con este bloque posterior del programa.

\section{Marcos de referencia y transformaciones geodésicas}

Un sistema de coordenadas indica cómo expresar posiciones; un marco de referencia materializa
ese sistema mediante convenciones, época y puntos o estaciones observadas. Cambiar de marco
puede exigir traslaciones, rotaciones, escala y correcciones espaciales.

\subsection{Transformaciones paramétricas}

Las transformaciones tridimensionales de semejanza suelen organizarse como
\[
  X_2=t+(1+\mu)RX_1,
\]
donde \(t\) es una traslación, \(\mu\) una pequeña variación de escala y \(R\) una rotación.
Las fórmulas de Molodensky permiten transformar coordenadas geodésicas considerando cambios
de datum y de elipsoide sin pasar necesariamente por una proyección plana.

\subsection{Transformaciones mediante grillas}

Una grilla NTv2 almacena correcciones \(\Delta\varphi,\Delta\lambda\) en nodos regulares. Para
un punto interior se interpolan las correcciones de los cuatro nodos vecinos, típicamente por
interpolación bilineal:
\[
\Delta(u,v)=
(1-u)(1-v)\Delta_{00}
+u(1-v)\Delta_{10}
+(1-u)v\Delta_{01}
+uv\Delta_{11}.
\]
Este enfoque representa deformaciones regionales que una transformación global de pocos
parámetros no puede modelar.

\section{Discretización, grillas y remuestreo}

Al pasar de coordenadas continuas a una grilla aparecen decisiones sobre resolución,
origen, orientación y regla de muestreo. Una transformación geométrica rara vez lleva los
centros de píxel exactamente a centros de otra grilla; por eso se requiere remuestreo.

\begin{center}
\begin{tabularx}{.9\textwidth}{>{\bfseries}l X}
\toprule
Método & Característica\\
\midrule
Vecino más próximo & conserva valores originales; apropiado para clases categóricas\\
Bilineal & interpola cuatro celdas; suaviza datos continuos\\
Bicúbico & usa un vecindario mayor; mayor suavidad y costo\\
\bottomrule
\end{tabularx}
\end{center}

El jacobiano de la transformación permite estudiar deformaciones locales y convertir áreas
y densidades entre sistemas de coordenadas.

\section{Sistemas de ecuaciones y aplicaciones}\label{sec:sistemas-ecuaciones}

El programa incluye sistemas de ecuaciones y ecuaciones de segundo orden. En el contexto de
coordenadas pueden aparecer al:
\begin{itemize}
  \item determinar intersecciones, puntos fijos y parámetros de transformación;
  \item ajustar coordenadas a partir de observaciones redundantes;
  \item describir curvas y superficies cuadráticas;
  \item formular ecuaciones diferenciales sobre superficies.
\end{itemize}

Este bloque deberá precisarse con los enunciados y ejemplos que presente la cátedra. Mientras tanto, los cuatro problemas del esqueleto 2025 (regresión, intersección de círculos, Procrustes, regresión total) están desarrollados, como anticipo, en la sección~\ref{sec:cuatro-problemas}.

\section{Conexión con la práctica de Agrimensura}

\begin{center}
\begin{tabularx}{\textwidth}{>{\bfseries}p{47mm} X}
\toprule
Concepto & Aplicación\\
\midrule
Coordenadas geodésicas \((\varphi,\lambda,h)\) & salida y procesamiento de observaciones GNSS\\
Base tangente y métrica & distancias y ángulos locales sobre el elipsoide\\
Radios \(M,N\) & variación geométrica con la latitud\\
Proyección conforme & coordenadas planas para cartografía y mensura\\
Transformaciones de marco & compatibilidad entre realizaciones geodésicas\\
Grillas de corrección & transformación regional de alta precisión\\
Remuestreo & ortomosaicos, modelos de elevación y datos ráster\\
Jacobiano & escala, deformación y propagación local\\
\bottomrule
\end{tabularx}
\end{center}

\begin{idea}[Estrategia de estudio acumulativa]
Cada nuevo tema debería incorporarse con la secuencia
\[
\text{definición}\rightarrow
\text{representación geométrica}\rightarrow
\text{fórmula}\rightarrow
\text{ejemplo numérico}\rightarrow
\text{implementación}.
\]
Para la implementación pueden usarse Python, NumPy, SciPy, SymPy, \texttt{pyproj},
GeographicLib y QGIS según el bloque.
\end{idea}


\section{Curvatura: de la curva plana a la forma cuadrática}\label{sec:curvatura}

{\small\vspace{1mm}\linkUnidad{12_curvatura}{Simulación interactiva: Unidad 16 — Curvatura como forma cuadrática}\par}
{\small\vspace{1mm}\linkAnimacion{05_elipsoide_euler_meusnier}{Animación: Euler y Meusnier sobre el elipsoide}\par}


\begin{atencion}[Anticipado]
Las clases 12, 13 y 14 del esqueleto (curvatura, elipsoide, trigonometría esférica) no se
dictaron en 2C2026. Las cuatro secciones que siguen se armaron antes de que la cátedra las
diera, a partir del resumen 2025 del profesor (una oración para la clase 12, sólo títulos
para las 13 y 14), de su apunte ``Geometría del Elipsoide y cálculos Geodésicos'' (guía
para clase, CC-BY-SA, figuras de Mariano Harguinteguy) y de bibliografía. Los valores
numéricos de control los calculamos con NumPy. Pendiente de validar.
\end{atencion}

El resumen 2025 de la clase 12 dice todo en una oración: ``la curvatura de una línea y la
idea de que la curvatura tiene que ser una forma cuadrática nos ayuda a entender la
curvatura del elipsoide y qué son las líneas geodésicas''. Son tres escalones y el orden
importa.

\paragraph{Curvatura de una curva plana.} El apunte del profesor la define como ``la tasa
de cambio del ángulo de la tangente a una curva con respecto al avance sobre dicha
curva'': si \(\alpha\) es el ángulo de la tangente con una dirección fija y \(s\) el arco,
\[
  \chi=\left|\frac{d\alpha}{ds}\right|,\qquad \text{circunferencia de radio } R:\ \chi=\frac1R .
\]
Hecho geométrico: en cada punto hay una circunferencia que besa a la curva (osculadora) y
la curvatura es \(1/R\) de esa circunferencia; en una recta, \(R=\infty\) y \(\chi=0\).
Para un gráfico \(y=f(x)\), con \(\tan\alpha=f'\) y \(ds=\sqrt{1+f'^2}\,dx\),
\[
  \chi=\frac{|f''|}{(1+f'^2)^{3/2}} ,
\]
y en un punto con tangente horizontal queda \(\chi=|f''|\): la segunda derivada \emph{es}
la curvatura. Parábola \(y=x^2/2\): en el origen \(\chi=1\), circunferencia osculadora de
radio \(1\) centrada en \((0,1)\); en \(x=1\), \(\chi=1/2^{3/2}\approx0{,}35\).

\paragraph{Prueba 1 de Clavijo: la curvatura normal es una forma cuadrática.} Sea una
superficie y \(p\) un punto. Se ponen ejes cartesianos con origen en \(p\) y \(z\) en la
dirección de la normal, de modo que cerca de \(p\) la superficie es \(z=f(x,y)\) con
\(f(0,0)=0\) y \(\nabla f(0,0)=0\) (plano tangente horizontal). Un plano \emph{normal}
contiene a la normal y a una dirección tangente \((\cos\alpha,\sin\alpha)\); corta a la
superficie en una curva plana cuya curvatura en \(p\) se llama curvatura normal en la
dirección \(\alpha\). Sobre esa curva, \(z(t)=f(t\cos\alpha,t\sin\alpha)\) tiene tangente
horizontal en \(t=0\), así que su curvatura es \(|z''(0)|\), y por la regla de la cadena
\[
  \boxed{\ \chi(\alpha)=\bigl|\,f_{xx}\cos^2\alpha+2f_{xy}\cos\alpha\sin\alpha+f_{yy}\sin^2\alpha\,\bigr|
  =\left|\begin{pmatrix}\cos\alpha&\sin\alpha\end{pmatrix}
  \begin{pmatrix}f_{xx}&f_{xy}\\f_{xy}&f_{yy}\end{pmatrix}
  \begin{pmatrix}\cos\alpha\\\sin\alpha\end{pmatrix}\right|\ }.
\]
Es la hessiana de la clase 9 evaluada en la dirección: un número que sale de un vector
pasando por una matriz simétrica, igual que \(\|v\|^2=v^Tv\) pero con otra matriz en el
medio. Eso es una forma cuadrática, y por eso la curvatura ``tiene que'' serlo. Con las
identidades del ángulo doble,
\[
  \chi(\alpha)=\left|\frac{f_{xx}+f_{yy}}2+\frac{f_{xx}-f_{yy}}2\cos2\alpha+f_{xy}\sin2\alpha\right|
  =\bigl|A+B\cos(2\alpha+C)\bigr| ,
\]
y ``independientemente de cuáles sean \(A\), \(B\) y \(C\), los máximos y mínimos se darán
a intervalos de \(\pi/2\)'': las direcciones de mayor y menor curvatura son perpendiculares.
Es lo mismo que decir que la hessiana, por simétrica, se diagonaliza con autovectores
ortogonales (sección~\ref{sec:cambio-base}); el resumen 2025 de la clase 10 ya lo dijo:
en las simétricas autovectores y vectores singulares coinciden.

\begin{atencion}[Errata confirmada en la Prueba 1 del apunte del elipsoide]
En la página 8 el apunte define \(L=f_{xx}\), \(N=f_{yy}\), \(M=f_{xy}\) y a renglón
seguido escribe \(\chi=|L\cos^2\alpha+2N\cos\alpha\sin\alpha+M\sin^2\alpha|\): el término
mixto lleva \(N\) y el de \(\sin^2\alpha\) lleva \(M\), al revés de las definiciones. La
fórmula es correcta si \(M\) es el mixto y \(N=f_{yy}\) (la notación habitual); la línea de
definiciones es la que está cambiada. Lo verificamos con \(f=\tfrac12x^2+2xy+\tfrac32y^2\)
(\(f_{xx}=1\), \(f_{xy}=2\), \(f_{yy}=3\)): la curvatura de la sección en \(\alpha=90^\circ\),
calculada por diferencias finitas, da \(3=f_{yy}\); la fórmula con las letras del apunte
da \(2\).
\end{atencion}

\begin{resultado}[Fórmula de Euler y radio medio]
Midiendo \(\alpha\) desde una dirección principal, con curvaturas principales
\(\chi_1\) y \(\chi_2\) (las de máxima y mínima curvatura, a \(90^\circ\) una de otra):
\[
  \chi(\alpha)=\chi_1\cos^2\alpha+\chi_2\sin^2\alpha
  =\frac{\cos^2\alpha}{R_1}+\frac{\sin^2\alpha}{R_2}.
\]
El apunte del profesor promedia el radio \(1/\chi(\alpha)\) sobre un cuadrante con la
sustitución \(u=\sqrt{R_2/R_1}\tan\alpha\), que convierte la integral en
\(\int_0^\infty du/(1+u^2)=\pi/2\), y obtiene el \emph{radio medio}
\[
  \boxed{\ R=\frac2\pi\int_0^{\pi/2}\frac{d\alpha}{\chi(\alpha)}=\sqrt{R_1R_2}\ }.
\]
En el elipsoide, \(R=\sqrt{MN}\): el radio de la esfera que localmente mejor lo reemplaza.
\end{resultado}

Ejemplo: \(z=x^2+xy+y^2\) en el origen tiene hessiana \(\begin{pmatrix}2&1\\1&2\end{pmatrix}\),
la misma matriz de la sección~\ref{sec:cambio-base}: curvaturas principales \(3\) y \(1\)
en las direcciones \((1,1)\) y \((1,-1)\); en la dirección del eje \(x\), a \(45^\circ\) de
la principal, Euler da \(3\cos^245^\circ+1\sin^245^\circ=2=f_{xx}\). La esfera
\(z=\sqrt{R^2-x^2-y^2}\) tiene en el polo hessiana \(-\tfrac1RI\): todas las direcciones
son principales, el caso que el apunte llama ``intuitivo, y no requería un análisis más
detallado, por el hecho de ser constante''.

\paragraph{Prueba 2: Meusnier.} El apunte remite la demostración a un apunte aparte
(Cristóbal de Lozada y Puga) y a una construcción en GeoGebra; acá sólo el enunciado y su
uso. Si un plano normal corta a la superficie con radio de curvatura \(\rho_\alpha\) y otro
plano, que comparte la misma recta tangente pero está inclinado un ángulo \(\theta\)
respecto del normal, la corta con radio \(\rho_\beta\), entonces
\[
  \boxed{\ \rho_\beta=\rho_\alpha\cos\theta\ }.
\]
Hecho geométrico: la circunferencia osculadora de la sección oblicua es la proyección de
la de la sección normal sobre el plano inclinado; proyectar acorta el radio por
\(\cos\theta\). Aplicación al paralelo: el paralelo de latitud \(\varphi\) es una
circunferencia de radio \(x'\) que comparte tangente (la dirección este-oeste) con la
sección normal perpendicular al meridiano, el \emph{primer vertical}; el plano del
paralelo está inclinado justo \(\varphi\) respecto de la normal. Por Meusnier,
\[
  x'=N\cos\varphi ,
\]
con \(N\) el radio del primer vertical. Es la deducción del apunte del profesor para
\(N\), y es lo que hace de \(N\) ``la distancia desde el punto hasta el eje polar medida a
lo largo de la normal''.

\paragraph{Por qué las direcciones principales del elipsoide son meridiano y primer
vertical.} Nuestra lectura del título de la clase 13, ``deducciones a partir de la simetría'':
el plano meridiano es un plano de simetría del elipsoide (reflejar en él lo deja igual).
Entonces la forma cuadrática de la curvatura en un punto es simétrica respecto de la
dirección del meridiano, y como sus extremos están a \(90^\circ\) uno de otro, tienen que
caer en el meridiano y en su perpendicular, sin calcular nada. Es lo que el apunte dice que
``mostraremos en forma sencilla''. La simetría de revolución agrega que nada métrico
depende de \(\lambda\): \(M\), \(N\) y \(K=1/(MN)\) son funciones de \(\varphi\) sola. Y en
los polos todas las direcciones son meridianos, así que ahí \(M=N\).

\section{Geometría del elipsoide: latitudes, radios y conversión}\label{sec:elipsoide}

{\small\vspace{1mm}\linkUnidad{13_elipsoide}{Simulación interactiva: Unidad 17 — Geometría del elipsoide}\par}
{\small\vspace{1mm}\linkAnimacion{05_elipsoide_euler_meusnier}{Animación: Euler y Meusnier sobre el elipsoide}\par}


\paragraph{La elipse meridiana y la latitud.} El elipsoide de revolución es la elipse
\(x'^2/a^2+z^2/b^2=1\) del plano meridiano, girada alrededor del eje \(z\): todos los
meridianos son la \emph{misma} elipse. Parámetros: aplastamiento \(f=(a-b)/a\), primera
excentricidad \(e^2=(a^2-b^2)/a^2=2f-f^2\), segunda \(e'^2=(a^2-b^2)/b^2\), con
\(b/a=\sqrt{1-e^2}=1/\sqrt{1+e'^2}\). La latitud geodésica \(\varphi\) es el ángulo entre
la \emph{normal} al elipsoide y el plano del ecuador. Hecho geométrico: la normal no pasa
por el centro, salvo en el ecuador y los polos; la plomada idealizada apunta ``más al polo''
que la recta al centro.

\begin{center}
\begin{tikzpicture}[scale=1.15]
  % elipse a=4, b=3; punto en phi=45: (3.2,1.8); normal corta eje x' en (1.4,0)
  \draw[axisline,-{Latex[length=2mm]}] (-0.4,0)--(4.9,0) node[right,font=\small] {\(x'\)};
  \draw[axisline,-{Latex[length=2mm]}] (0,-0.3)--(0,3.5) node[above,font=\small] {\(z\)};
  \draw[auxline] (0,0) circle (4);
  \draw[teal,line width=1.1pt] (0,0) ellipse (4 and 3);
  \coordinate (p) at (3.2,1.8);
  \coordinate (q) at (3.2,2.4);
  \coordinate (n0) at (1.4,0);
  \draw[auxline] (3.2,0)--(q);
  \draw[rust,line width=1pt] (n0)--($(p)+(0.55,0.55)$) node[above right,font=\small,labelback] {normal};
  \draw[auxline] ($(p)+(-1.2,1.2)$)--($(p)+(1.0,-1.0)$);
  \draw[tealdark,line width=.9pt] (0,0)--(p);
  \draw[gold,line width=.9pt] (0,0)--(q);
  \fill[ink] (p) circle (1.8pt) node[above left,font=\small,labelback] {\(P=(x',z)\)};
  \fill[gold] (q) circle (1.6pt) node[above,font=\small,labelback] {\(Q\)};
  \draw[rust] (2.0,0) arc (0:45:0.6) node[midway,right,font=\small,xshift=1pt] {\(\varphi\)};
  \draw[tealdark] (1.25,0) arc (0:29.4:1.25) node[midway,right,font=\small,xshift=2pt] {\(\psi\)};
  \draw[gold] (0.8,0) arc (0:36.9:0.8) node[midway,above left,font=\small,xshift=3pt,yshift=-2pt] {\(\beta\)};
  \node[text=muted,font=\small] at (3.9,-0.45) {elipse meridiana};
  \node[text=muted,font=\small] at (1.7,3.55) {circunferencia auxiliar de radio \(a\)};
\end{tikzpicture}
\end{center}

La recta tangente a la elipse en \(P\) es perpendicular a la normal, así que su pendiente
es \(\tan(\varphi+\pi/2)=-\cot\varphi\). Derivando la elipse,
\(dz/dx'=-(b^2/a^2)(x'/z)\), e igualando:
\[
  \frac{b^2}{a^2}\frac{x'}{z}=\cot\varphi
  \quad\Longrightarrow\quad
  z=\frac{b^2}{a^2}\,x'\tan\varphi,\qquad
  z^2=\frac{\sin^2\varphi}{\cos^2\varphi}\,\frac{b^4}{a^4}\,x'^2 ,
\]
que es la ecuación (2) del apunte del profesor. Reemplazando en la elipse y despejando:
\[
  \boxed{\ x'=\frac{a\cos\varphi}{\sqrt{1-e^2\sin^2\varphi}}=N\cos\varphi\ },\qquad
  \boxed{\ z=\frac{b\sin\varphi}{\sqrt{1+e'^2\cos^2\varphi}}
  =\frac{a(1-e^2)\sin\varphi}{\sqrt{1-e^2\sin^2\varphi}}=N(1-e^2)\sin\varphi\ }.
\]
La igualdad de las dos formas de \(z\) sale de
\(1+e'^2\cos^2\varphi=(1-e^2\sin^2\varphi)/(1-e^2)\). Con la longitud \(\lambda\), medida
desde el meridiano origen (el que contiene al eje \(x\)) hasta el plano meridiano del punto:
\[
  x=x'\cos\lambda,\qquad y=x'\sin\lambda,\qquad z=N(1-e^2)\sin\varphi .
\]
Es la parametrización \(r(\varphi,\lambda)\) del elipsoide, la que usa la
sección~\ref{sec:curvilineas}. Un punto con altura \(h\) sobre el elipsoide se corre \(h\)
a lo largo de la normal: \(N\to N+h\) en \(x,y\) y \(N(1-e^2)\to N(1-e^2)+h\) en \(z\).

\begin{atencion}[Dos erratas confirmadas en la conversión del apunte del elipsoide]
\begin{itemize}
  \item En la página 15, la expresión de \(z\) ``en términos de \(e^2\)'' lleva
  \(\sqrt{1-e^2\cos^2\varphi}\) en el denominador; el correcto es
  \(\sqrt{1-e^2\sin^2\varphi}\), el mismo denominador que \(x'\) y que \(N\). Lo verificamos
  de dos maneras: algebraicamente, desde la forma en \(e'^2\) que el propio apunte da dos
  líneas antes; y numéricamente en GRS80 para \(\varphi=-34{,}6^\circ\): la fórmula correcta
  da \(z=-3\,601\,428{,}7\) m y la errónea \(-3\,605\,726{,}3\) m, una diferencia de
  \(4{,}3\) km. En el polo la errónea da \(b^2/a=6\,335\,439\) m en vez de \(b\).
  \item En la misma página, ``\(x=x'\sin\lambda\), \(y=x'\cos\lambda\)'' invierte seno y
  coseno respecto de la convención habitual (eje \(x\) hacia el meridiano origen) y de
  la propia figura 3 del apunte, que marca \(\lambda\) desde el eje \(x\) hasta la
  dirección \(x'\). Con esa figura, \(x=x'\cos\lambda\) e \(y=x'\sin\lambda\).
\end{itemize}
\end{atencion}

\paragraph{Los radios \(M\) y \(N\).} El radio del primer vertical salió por Meusnier:
\(N=a/\sqrt{1-e^2\sin^2\varphi}\). El radio de curvatura del meridiano, \(M\), es el de
la elipse meridiana en \(P\); el apunte del profesor lo deduce aplicando
\(\chi=|z''|/(1+z'^2)^{3/2}\) a \(z(x')\), con \(z''=-\tfrac1z\tfrac{b^2}{a^2}(1+\tfrac{a^2}{b^2}z'^2)\)
y reemplazando \(z'=-\cot\varphi\) y el \(z(\varphi)\) recién hallado:
\[
  \boxed{\ M=\frac{a(1-e^2)}{(1-e^2\sin^2\varphi)^{3/2}}\ },\qquad
  \boxed{\ N=\frac{a}{\sqrt{1-e^2\sin^2\varphi}}\ },\qquad
  M\le N,\ \text{con igualdad sólo en los polos}.
\]
Controles: en el ecuador \(M=a(1-e^2)=b^2/a\) (el radio osculador de la elipse en el
extremo del eje mayor, el mínimo) y \(N=a\); en el polo \(M=N=a/\sqrt{1-e^2}=a^2/b\), el
máximo. El elipsoide es más curvo a lo largo del meridiano que a lo ancho. Además,
\(\|r_\varphi\|=M\) y \(\|r_\lambda\|=x'=N\cos\varphi\), de modo que la primera forma
fundamental de la sección~\ref{sec:metrica} es
\[
  ds^2=M^2\,d\varphi^2+N^2\cos^2\varphi\,d\lambda^2 ,
\]
y la esfera, \(ds^2=R^2d\varphi^2+R^2\cos^2\varphi\,d\lambda^2\), es el caso \(M=N=R\).

\paragraph{Las tres latitudes.} Además de la geodésica \(\varphi\) (normal), la
\emph{geocéntrica} \(\psi\), ángulo de la recta al centro, y la \emph{reducida} \(\beta\),
la del punto \(Q\) de la circunferencia auxiliar de radio \(a\) que está sobre la misma
vertical que \(P\) (\(x'=a\cos\beta\), \(z=b\sin\beta\)):
\[
  \tan\psi=\frac{z}{x'}=(1-e^2)\tan\varphi,\qquad
  \tan\beta=\frac ab\,\frac{z}{x'}=\sqrt{1-e^2}\,\tan\varphi,\qquad
  |\psi|\le|\beta|\le|\varphi| .
\]
La geocéntrica es siempre menor en módulo que la geodésica, y la diferencia es máxima
cerca de \(45^\circ\).

\begin{resultado}[Valores de control en GRS80 (\(a=6\,378\,137\) m, \(f=1/298{,}257222101\))]
\(e^2=0{,}00669438\), \(b=6\,356\,752{,}314\) m. Calculados con NumPy; la fila de
\(-34{,}6^\circ\) es la latitud de Buenos Aires.

{\small
\begin{center}
\begin{tabular}{l rrrr}
\toprule
 & \(\varphi=0^\circ\) & \(\varphi=45^\circ\) & \(\varphi=-34{,}6^\circ\) & \(\varphi=90^\circ\)\\
\midrule
\(M\) (m) & 6\,335\,439{,}327 & 6\,367\,381{,}816 & 6\,356\,008{,}141 & 6\,399\,593{,}626\\
\(N\) (m) & 6\,378\,137{,}000 & 6\,388\,838{,}290 & 6\,385\,032{,}023 & 6\,399\,593{,}626\\
\(x'=N\cos\varphi\) (m) & 6\,378\,137{,}000 & 4\,517\,590{,}879 & 5\,255\,752{,}072 & 0\\
\(z=N(1-e^2)\sin\varphi\) (m) & 0 & 4\,487\,348{,}409 & \(-3\,601\,428{,}679\) & 6\,356\,752{,}314\\
\(\psi\) (geocéntrica) & \(0^\circ\) & \(44{,}807577^\circ\) & \(-34{,}420332^\circ\) & \(90^\circ\)\\
\(\beta\) (reducida) & \(0^\circ\) & \(44{,}903788^\circ\) & \(-34{,}510112^\circ\) & \(90^\circ\)\\
\(|\varphi-\psi|\) & \(0'\) & \(11{,}545'\) & \(10{,}780'\) & \(0'\)\\
\(R=\sqrt{MN}\) (m) & 6\,356\,752{,}314 & 6\,378\,101{,}030 & 6\,370\,503{,}553 & 6\,399\,593{,}626\\
\bottomrule
\end{tabular}
\end{center}
}

Chequeo: en cada columna \(x'^2/a^2+z^2/b^2=1\) con doce decimales. Un mojón en
\(\varphi=-34{,}6^\circ\), \(\lambda=-58{,}4^\circ\), \(h=0\) tiene
\((X,Y,Z)=(2\,753\,940{,}011,\ -4\,476\,465{,}599,\ -3\,601\,428{,}679)\) m. La diferencia
\(N-M\) en el ecuador es \(ae^2=42\,698\) m.
\end{resultado}

\section{Arco de meridiano}\label{sec:arco-meridiano}

\textbf{Campo.} Cuánto mide, en metros, un tramo de meridiano entre dos latitudes: es lo
que necesita cualquier proyección que respete distancias a lo largo del meridiano (el
``puntapié inicial'' del apunte del profesor hacia Mercator transversa y Gauss--Krüger).
Como \(\|r_\varphi\|=M\),
\[
  S=\int_{\varphi_0}^{\varphi_1}M\,d\varphi
  =a(1-e^2)\int_{\varphi_0}^{\varphi_1}\bigl(1-e^2\sin^2\varphi\bigr)^{-3/2}d\varphi ,
\]
integral que no tiene primitiva elemental. El camino del apunte: serie binomial del
integrando, que deja potencias pares del seno; pasaje de esas potencias a cosenos de
ángulos múltiples, que sí se integran.

\paragraph{Serie binomial.} \((1+x)^\alpha=\sum_k\binom{\alpha}{k}x^k\) con
\(\alpha=-\tfrac32\) y \(x=-e^2\sin^2\varphi\). Los signos se cancelan (el coeficiente
binomial de exponente negativo alterna de signo, y \(x^k\) también) y quedan todos
positivos:
\[
  \bigl(1-e^2\sin^2\varphi\bigr)^{-3/2}
  =\sum_{n\ge0}\frac{(2n+1)!}{(2^n\,n!)^2}\,e^{2n}\sin^{2n}\varphi
  =1+\tfrac32e^2\sin^2\varphi+\tfrac{15}8e^4\sin^4\varphi+\tfrac{35}{16}e^6\sin^6\varphi
  +\tfrac{315}{128}e^8\sin^8\varphi+\cdots
\]
(verificado término a término: \(\tfrac32\cdot\tfrac52/2!=\tfrac{15}8\),
\(\tfrac32\cdot\tfrac52\cdot\tfrac72/3!=\tfrac{35}{16}\), etc.).

\paragraph{Potencias del seno como cosenos de ángulos múltiples.}
\[
  \sin^2\varphi=\frac{1-\cos2\varphi}2,\qquad
  \sin^4\varphi=\frac{3-4\cos2\varphi+\cos4\varphi}8,
\]
\[
  \sin^6\varphi=\frac{10-15\cos2\varphi+6\cos4\varphi-\cos6\varphi}{32},\qquad
  \sin^8\varphi=\frac{35-56\cos2\varphi+28\cos4\varphi-8\cos6\varphi+\cos8\varphi}{128}.
\]
Agrupando por \(\cos2k\varphi\), el integrando queda
\(A+B\cos2\varphi+C\cos4\varphi+D\cos6\varphi+E\cos8\varphi+\cdots\), y el apunte del
profesor escribe cada coeficiente con una sola plantilla, ``para minimizar los errores a la
hora de programar las fórmulas en una planilla'':
\[
  \text{coeficiente}=\sum_{n}\frac{x_n}{2^{\,2n-1}}\,\frac{(2n+1)!}{(2^n\,n!)^2}\,e^{2n},
\]
con los \(x_n\) de la tabla (transcrita del apunte y verificada contra las identidades de
arriba: cada fila es la columna correspondiente de los numeradores, y \(2^{2n-1}\) es el
denominador \(2,8,32,128\); para \(n=0\) vale \(x_0=\tfrac12\) porque \(2^{-1}\) lo
duplica).

{\small
\begin{center}
\begin{tabular}{l rrrrr}
\toprule
Término & \(x_0\) & \(x_1\) & \(x_2\) & \(x_3\) & \(x_4\)\\
\midrule
\(A\) & \(\tfrac12\) & 1 & 3 & 10 & 35\\
\(B\) & 0 & \(-1\) & \(-4\) & \(-15\) & \(-56\)\\
\(C\) & 0 & 0 & 1 & 6 & 28\\
\(D\) & 0 & 0 & 0 & \(-1\) & \(-8\)\\
\(E\) & 0 & 0 & 0 & 0 & 1\\
\bottomrule
\end{tabular}
\end{center}
}

\begin{resultado}[Longitud del arco de meridiano]
\[
  S=a(1-e^2)\Bigl[A\,\varphi+\tfrac B2\sin2\varphi+\tfrac C4\sin4\varphi
  +\tfrac D6\sin6\varphi+\tfrac E8\sin8\varphi+\cdots\Bigr]_{\varphi_0}^{\varphi_1}
  \qquad(\varphi\ \text{en radianes}).
\]
Los coeficientes dependen sólo de \(e^2\): se calculan una vez por elipsoide. En GRS80,
\(A=1{,}0050525018\), \(B=-5{,}0631086\cdot10^{-3}\), \(C=1{,}0627582\cdot10^{-5}\),
\(D=-2{,}08172\cdot10^{-8}\), \(E=3{,}861\cdot10^{-11}\). Controles contra la integral
numérica (\texttt{scipy.integrate.quad}): del ecuador a \(45^\circ\),
\(S=4\,984\,944{,}378\) m; del ecuador a \(-34{,}6^\circ\), \(3\,830\,218{,}134\) m; el
cuadrante de meridiano, \(10\,001\,965{,}729\) m. En los tres casos la serie con cinco
términos coincide con la integral en menos de \(0{,}1\) mm. El ejercicio que deja el
apunte (graficar el decrecimiento de los \(x_n\) y estimar en metros el aporte de cada
término) sale de la misma cuenta: de \(0\) a \(30^\circ\) el término \(A\) aporta
\(3\,333\,988{,}58\) m, el \(B\) \(-13\,889{,}76\) m, el \(C\) \(14{,}58\) m, el \(E\)
\(-2{,}6\cdot10^{-5}\) m (el \(D\) se anula ahí porque \(\sin180^\circ=0\)).
\end{resultado}

\begin{atencion}[Errata confirmada en el arco de meridiano del apunte del elipsoide]
En la página 18 el apunte escribe
\(\sin^6\varphi=(10-15\cos2\varphi+6\cos4\varphi+\cos6\varphi)/32\); el término
\(\cos6\varphi\) lleva signo \emph{menos}. Se ve en \(\varphi=0\): con el signo más la
expresión vale \(2/32\neq0\). Es coherente con la propia tabla del apunte, donde
\(x_3=-1\) en la fila \(D\). Y en la primitiva de la página 19 falta multiplicar \(A\) por
\(\varphi\): dice ``\(A+2^{-1}B\sin2\varphi+\cdots\)'' donde corresponde
``\(A\varphi+2^{-1}B\sin2\varphi+\cdots\)''. Con esas dos correcciones la fórmula reproduce
la integral numérica, como muestran los controles de arriba.
\end{atencion}

\paragraph{Los cinco problemas de proyección del cierre del apunte.} El apunte del
profesor termina con un ``puntapié inicial'': con la conversión, los módulos
\(\|r_\varphi\|=M\), \(\|r_\lambda\|=N\cos\varphi\) y los arcos, ``¿cómo plantearían las
condiciones para desarrollar una proyección sobre el elipsoide?'', y propone bosquejar
cinco (contamos cinco en la página 20). Los planteamos acá; en las tres cilíndricas la cuenta la
verificamos contra la integral numérica y contra \texttt{pyproj} (proyecciones
\texttt{merc}, \texttt{eqc} y \texttt{cea} sobre GRS80), que coinciden al milímetro. En
las tres, \(x=a\lambda\) (el ecuador se desarrolla sin deformación) y lo que cambia es
\(y(\varphi)\).
\begin{enumerate}
  \item \textbf{Mercator (conforme).} Condición: la escala sea la misma en meridiano y
  paralelo,
  \[
    \frac{dy}{M\,d\varphi}=\frac{dx}{N\cos\varphi\,d\lambda}
    \quad\Longrightarrow\quad dy=\frac{a\,M}{N\cos\varphi}\,d\varphi ,
  \]
  e integrando:
  \[
    y=a\Bigl[\ln\tan\bigl(\tfrac\pi4+\tfrac\varphi2\bigr)
    -\tfrac e2\ln\frac{1+e\sin\varphi}{1-e\sin\varphi}\Bigr].
  \]
  Control: \(y(45^\circ)=5\,591\,295{,}918\) m (la esférica, sin el segundo término,
  daría \(5\,621\,521{,}486\)); \(y(-34{,}6^\circ)=-4\,085\,391{,}659\) m. La escala en el
  paralelo es \(a/(N\cos\varphi)\): \(1{,}4118\) en \(45^\circ\), \(1{,}2136\) en Buenos
  Aires.
  \item \textbf{Cilíndrica equidistante.} Condición: los meridianos a escala \(1\),
  \(dy=M\,d\varphi\), o sea \(y=S(\varphi)\), el arco de meridiano de arriba:
  \(y(45^\circ)=4\,984\,944{,}378\) m, \(y(-34{,}6^\circ)=-3\,830\,218{,}134\) m.
  \item \textbf{Cilíndrica equivalente.} Condición: el área se conserve,
  \[
    dx\,dy=MN\cos\varphi\,d\varphi\,d\lambda\quad\Longrightarrow\quad
    dy=\frac{MN\cos\varphi}{a}\,d\varphi .
  \]
  Integrando, \(y=\tfrac a2\,q(\varphi)\) con
  \[
    q(\varphi)=(1-e^2)\Bigl[\frac{\sin\varphi}{1-e^2\sin^2\varphi}
    -\frac1{2e}\ln\frac{1-e\sin\varphi}{1+e\sin\varphi}\Bigr]:
  \]
  \(y(45^\circ)=4\,489\,858{,}887\) m, \(y(-34{,}6^\circ)=-3\,602\,726{,}723\) m.
  \item \textbf{Elipsoide a esfera sin modificar las áreas.} Condición: una esfera de
  radio \(R_q\) con la misma área total, \(4\pi R_q^2=2\pi a^2q(\pi/2)\), es decir
  \(R_q=a\sqrt{q(\pi/2)/2}=6\,371\,007{,}181\) m (control: el área del elipsoide GRS80
  es \(510\,065\,621{,}7\) km\(^2\) por la fórmula exacta
  \(2\pi a^2+\pi\frac{b^2}e\ln\frac{1+e}{1-e}\), y \(4\pi R_q^2\) da lo mismo); y a cada
  paralelo \(\varphi\) le corresponde en la esfera la latitud \(\beta_q\) con
  \(\sin\beta_q=q(\varphi)/q(\pi/2)\) (latitud autálica): \(44{,}8717^\circ\) para
  \(45^\circ\), \(-34{,}4801^\circ\) para Buenos Aires. Lo demás (las longitudes) no cambia.
  \item \textbf{Transversa de Mercator.} Condiciones: conforme (misma escala en todas las
  direcciones en cada punto), y el meridiano central sin deformación, con \(y=S(\varphi)\)
  sobre él. Es la que usa Gauss--Krüger y la cartografía oficial argentina. No la resolvemos:
  sobre el elipsoide no tiene fórmula cerrada y se desarrolla en serie de potencias de la
  diferencia de longitud al meridiano central; queda como bosquejo, igual que en el apunte.
\end{enumerate}

\section{Polígonos geodésicos y trigonometría esférica}\label{sec:trigonometria-esferica}

{\small\vspace{1mm}\linkUnidad{14_trigonometria_esferica}{Simulación interactiva: Unidad 18 — Trigonometría esférica}\par}


\textbf{Campo.} El apunte del profesor arranca con un observador que construye ``una
poligonal cerrada sobre la tierra, verificando si el error obtenido al regresar al punto de
origen es compatible con la precisión de los instrumentos''. Entre vértices lejanos los
lados son geodésicas (las ``rectas'' de la superficie: la curva más corta entre dos puntos,
la que no se tuerce dentro de la superficie) y la suma de los ángulos de la poligonal no da
lo que daría en el plano. Sobre la esfera hay fórmulas cerradas; sobre el elipsoide no, y
ahí entra el álgebra simbólica (el ``Máxima para el elipsoide'' del título de la clase 14).

\paragraph{Hecho geométrico: un triángulo esférico son tres planos por el centro.} Sean
\(A,B,C\) tres vectores unitarios desde el centro de la esfera. Los lados \(a,b,c\)
(opuestos a cada vértice) son arcos de círculo máximo, es decir \emph{ángulos} entre
vectores: \(\cos a=B\cdot C\), \(\cos b=A\cdot C\), \(\cos c=A\cdot B\)
(sección~\ref{sec:interno}). El ángulo en \(A\) es el diedro entre los planos \(OAB\) y
\(OAC\), es decir el ángulo entre sus normales \(A\times B\) y \(A\times C\)
(sección~\ref{sec:vectorial}). Nada mide en metros: todo son ángulos, y la esfera de radio
\(R\) sólo multiplica los lados por \(R\).

\begin{resultado}[Teoremas del coseno y del seno esféricos]
\[
  \cos a=\cos b\cos c+\sin b\sin c\cos A,\qquad
  \frac{\sin a}{\sin A}=\frac{\sin b}{\sin B}=\frac{\sin c}{\sin C}.
\]
Deducción del coseno con un solo producto interno: ponemos \(A\) en el polo, \(A=(0,0,1)\);
\(B\) a distancia \(c\) sobre el meridiano \(\lambda=0\), \(B=(\sin c,0,\cos c)\); \(C\)
a distancia \(b\) sobre el meridiano \(\lambda=A\) (el ángulo en el polo es la diferencia
de longitudes), \(C=(\sin b\cos A,\sin b\sin A,\cos b)\). Entonces
\(\cos a=B\cdot C=\sin b\sin c\cos A+\cos b\cos c\). Sin coordenadas:
\((A\times B)\cdot(A\times C)=(A\cdot A)(B\cdot C)-(A\cdot B)(A\cdot C)=\cos a-\cos b\cos c\),
y \(\|A\times B\|=\sin c\), \(\|A\times C\|=\sin b\). El del seno sale del volumen
\(|\det[A\,B\,C]|\), el mismo determinante de la sección~\ref{sec:determinante}:
\(\|(A\times B)\times(A\times C)\|=|\det[A\,B\,C]|=\sin b\sin c\sin A\), y por simetría
también \(=\sin a\sin c\sin B=\sin a\sin b\sin C\). Para lados chicos
(\(\sin b\approx b\), \(\cos b\approx1-b^2/2\)) el coseno esférico devuelve
\(a^2=b^2+c^2-2bc\cos A\): una Tierra lo bastante chica es plana.
\end{resultado}

\paragraph{Exceso esférico.} En el plano los ángulos de un triángulo suman \(\pi\); en la
esfera suman más, y el sobrante es el área (Girard):
\[
  \boxed{\ \varepsilon=A+B+C-\pi=\frac{\text{Área}}{R^2}\ },\qquad
  \text{polígono de } n \text{ lados geodésicos: }\ \sum\text{ángulos}-(n-2)\pi=\frac{\text{Área}}{R^2}.
\]
Hecho geométrico: el triángulo ``abomba'' y los ángulos se abren; cuanto más grande, más
se abren. Control: el octante \(A=(0,0,1)\), \(B=(1,0,0)\), \(C=(0,1,0)\) tiene
\(b=c=90^\circ\), \(A=90^\circ\), \(\cos a=0\) y los tres ángulos rectos: suman
\(270^\circ\), exceso \(\pi/2\), área \(\pi R^2/2\), un octavo de la esfera. Un triángulo
equilátero de \(1\) km de lado sobre \(R=6371\) km tiene área \(0{,}433\) km\(^2\) y exceso
\(1{,}07\cdot10^{-8}\) rad \(\approx0{,}002''\); con lados de \(100\) km, área
\(4330\) km\(^2\) y exceso \(1{,}07\cdot10^{-4}\) rad \(\approx22''\). En una poligonal de
agrimensura el ``error angular de cierre'' tiene una parte que no es error sino exceso
esférico, y hay que separarla antes de compensar (nota nuestra, no del profesor: el teorema
de Legendre resuelve un triángulo esférico chico como plano restando \(\varepsilon/3\) a
cada ángulo).

\begin{enunciado}[Problema resuelto: exceso esférico de un triángulo entre tres ciudades]
Sobre la esfera de radio \(R=6371\) km, el triángulo con vértices en Buenos Aires
\((-34{,}6^\circ,-58{,}4^\circ)\), Córdoba \((-31{,}4^\circ,-64{,}2^\circ)\) y Mendoza
\((-32{,}9^\circ,-68{,}8^\circ)\). Hallar lados, ángulos, exceso esférico y área, y comparar
el ángulo en Buenos Aires con el que daría un triángulo plano de los mismos lados.
\end{enunciado}

Cada vértice como vector unitario,
\((\cos\varphi\cos\lambda,\ \cos\varphi\sin\lambda,\ \sin\varphi)\). Lados por producto
interno, \(\cos a=B\cdot C\), etc.: \(a\) (Córdoba--Mendoza) \(=4{,}1730^\circ=464{,}02\) km,
\(b\) (Buenos Aires--Mendoza) \(=8{,}8085^\circ=979{,}46\) km, \(c\) (Buenos
Aires--Córdoba) \(=5{,}8208^\circ=647{,}24\) km. Ángulos por el ángulo entre las normales
\(A\times B\) y \(A\times C\): en Buenos Aires \(23{,}5351^\circ\), en Córdoba
\(122{,}8281^\circ\), en Mendoza \(33{,}8154^\circ\). Suma \(180{,}1785^\circ\): el exceso
es \(\varepsilon=0{,}1785^\circ=642{,}7''=3{,}1159\cdot10^{-3}\) rad, y el área
\(\varepsilon R^2=126\,475\) km\(^2\). Controles: la fórmula de L'Huilier,
\(\tan\frac\varepsilon4=\sqrt{\tan\frac s2\tan\frac{s-a}2\tan\frac{s-b}2\tan\frac{s-c}2}\)
con \(s\) el semiperímetro angular, da el mismo \(\varepsilon\) con doce decimales; el
teorema del coseno cierra, \(\cos a=0{,}997349\) por los dos lados; y los tres cocientes
\(\sin a/\sin A\) valen \(0{,}18223\). El triángulo plano con los mismos lados tiene en
Buenos Aires \(23{,}4755^\circ\); el esférico menos un tercio del exceso,
\(23{,}5351^\circ-0{,}0595^\circ=23{,}4756^\circ\): Legendre funciona con una diferencia de
\(0{,}02''\). En una poligonal entre estas tres ciudades, \(10'43''\) de la suma angular son
geometría de la esfera, no error de medición.

\paragraph{La esfera local del elipsoide.} Sobre el elipsoide los lados de la poligonal
son geodésicas que no son círculos (los meridianos sí lo son; los paralelos no, salvo el
ecuador) y no hay teorema del coseno cerrado: los problemas directo (punto, acimut y
distancia \(\to\) punto) e inverso (dos puntos \(\to\) distancia y acimutes) se resuelven
con series como la del arco de meridiano, o numéricamente. Para una zona de trabajo
alcanza con reemplazar el elipsoide por la esfera de radio medio
\(R=\sqrt{MN}\) de la sección~\ref{sec:curvatura}: en \(\varphi=-34{,}6^\circ\),
\(R=6\,370\,503{,}6\) m, y el exceso del triángulo de \(100\) km recalculado con ese radio
es \(1{,}067\cdot10^{-4}\) rad, los mismos \(22''\). Qué se conserva: al rotar la esfera,
lados, ángulos y área del triángulo no cambian (el producto interno es invariante por
rotaciones, primera clase); lo que no se conserva al pasar de esfera a elipsoide es la
simetría total, y por eso se pierden las fórmulas cerradas. La unidad
\texttt{14\_trigonometria\_esferica} (apéndice~\ref{ap:notebooks}) hace la cuenta con
tres vértices movibles.

\section{Transformaciones conformes: del número complejo y el Jacobiano a la proyección del elipsoide}\label{sec:transformaciones-conformes}

{\small\vspace{1mm}\linkUnidad{15_transformaciones_conformes}{Simulación interactiva: Unidad 19 — Transformaciones conformes: números complejos, Jacobiano y del elipsoide al plano cartográfico}\par}
{\small\linkAnimacion{06_conforme_tissot_mercator}{Animación: indicatriz de Tissot, equidistante vs. Mercator}\par}


\begin{atencion}[Desarrollo complementario]
Esta sección no fue dictada por la cátedra en 2C2026. La armamos nosotros para cerrar el
arco de la materia: el número complejo de la clase 1 (sección~\ref{sec:rot-complejos}),
las similaridades de la clase 3 (secciones~\ref{sec:escalar} y~\ref{cap:08}), la SVD de la
clase 7 (sección~\ref{sec:svd}), el Jacobiano y la primera forma de la clase 9
(sección~\ref{sec:metrica}) y el elipsoide de la clase 13
(sección~\ref{sec:elipsoide}) son, vistos desde acá, cuatro maneras de decir lo mismo. Y
con ese lenguaje resolvemos las preguntas con que termina el apunte del elipsoide del
profesor, incluida la transversa de Mercator que en la sección~\ref{sec:arco-meridiano}
dejamos planteada. Pendiente de validar por la cátedra.
\end{atencion}

\paragraph{El problema cartográfico.} En la sección~\ref{sec:curvatura} vimos que la
curvatura gaussiana del elipsoide es \(K=1/(MN)>0\) en todos los puntos; la del plano es
\(K=0\). El \emph{theorema egregium} de Gauss dice que \(K\) se puede calcular con la
primera forma fundamental sola (con \(E,F,G\) y sus derivadas), es decir, midiendo
longitudes \emph{dentro} de la superficie, sin salir a \(\R^3\). La consecuencia es
inmediata: si existiera una aplicación del elipsoide al plano que conservara todas las
longitudes (una isometría local), las dos superficies tendrían la misma \(K\), y
\(1/(MN)\neq0\). No existe. Toda carta, y todo sistema de coordenadas planas
(Gauss--Krüger en las fajas de POSGAR), deforma \emph{algo}: longitudes, áreas, ángulos o
una mezcla. Hecho geométrico: una cáscara de naranja no se aplana sin rajarse.

Entonces elegimos qué conservar. Un agrimensor no negocia el ángulo: los \(90^\circ\)
medidos con la estación total entre dos alineaciones tienen que cortarse a \(90^\circ\) en
el plano de la mensura, y un acimut medido en el terreno tiene que ser el acimut del
plano más una corrección que dependa del punto, no de la dirección. Eso es una
\emph{transformación conforme}: en cada punto conserva los ángulos entre curvas, con su
orientación (no refleja). Lo que sí cambia de punto a punto es la escala, y la pagamos con
un factor \(k(P)\) que hay que conocer (la ``escala de la proyección'').

\paragraph{Qué es una transformación en un punto.} Una aplicación \(P\mapsto P'\) entre
superficies (o entre planos) es, cerca de cada punto, su jacobiana \(J(P)\): la matriz
\(2\times2\) que lleva un pasito \((du,dv)\) del terreno al pasito correspondiente del mapa
(sección~\ref{sec:metrica}). Para que ``conservar ángulos'' tenga sentido métrico, medimos
los pasitos del terreno con la primera forma del terreno y los del mapa con la del mapa; en
lo que sigue suponemos que trabajamos en bases ortonormales de los dos planos tangentes
(sobre el elipsoide, \(r_\varphi/M\) y \(r_\lambda/(N\cos\varphi)\); en el plano, los
versores), de modo que \(J\) es una matriz real \(2\times2\) común y corriente y la
pregunta ``¿conserva ángulos?'' es la de la clase 3.

\begin{resultado}[Las cuatro caras de una transformación conforme]
Sea \(J=J(P)\) la jacobiana en un punto, invertible. Son equivalentes:
\begin{enumerate}[label=(\roman*)]
  \item \textbf{Similaridad directa.} \(J=\rho\,R(\theta)\) con \(\rho>0\): una rotación
  seguida de una homotecia, la similaridad directa con punto fijo de la clase 3
  (sección~\ref{sec:escalar}).
  \item \textbf{Valores singulares iguales.} \(\sigma_1=\sigma_2=k\) en la SVD de \(J\)
  (sección~\ref{sec:svd}): la indicatriz de Tissot (la imagen de la circunferencia unidad
  del terreno) es una circunferencia de radio \(k\), no una elipse.
  \item \textbf{Métrica proporcional.} \(J^TJ=k^2I\): en la primera forma inducida,
  \(F=0\) y \(E=G=k^2\), es decir \(ds^2_{\text{mapa}}=k^2\,ds^2_{\text{terreno}}\).
  \item \textbf{Número complejo.} \(J=\begin{pmatrix}a&-b\\b&a\end{pmatrix}\), la matriz
  de multiplicar por \(a+ib\) (sección~\ref{sec:rot-complejos}). Si la transformación es
  \((u,v)\mapsto(x,y)\), esto son las ecuaciones de Cauchy--Riemann
  \[
    x_u=y_v,\qquad x_v=-y_u ,
  \]
  y \(J\) es multiplicar por \(f'(z)=x_u+i\,y_u\), con \(z=u+iv\) y \(f=x+iy\): toda
  función analítica con \(f'(z)\neq0\) es conforme, y recíprocamente.
\end{enumerate}
En todas, \(k=\rho=|f'(z)|=\sqrt{\det J}\) es la escala local y \(\theta=\arg f'(z)\) la
rotación local (la ``convergencia de meridianos'' en las proyecciones).
\end{resultado}

\emph{Demostración.} Son tres equivalencias con el centro en (i).

(i)\(\Leftrightarrow\)(iv), por la forma de las matrices. Una rotación por homotecia es
\[
  \rho\,R(\theta)=\begin{pmatrix}\rho\cos\theta&-\rho\sin\theta\\ \rho\sin\theta&\rho\cos\theta\end{pmatrix}
  =\begin{pmatrix}a&-b\\b&a\end{pmatrix}\quad\text{con } a=\rho\cos\theta,\ b=\rho\sin\theta ,
\]
y al revés: toda matriz \(\begin{pmatrix}a&-b\\b&a\end{pmatrix}\) con \((a,b)\neq(0,0)\)
se escribe así tomando \(\rho=\sqrt{a^2+b^2}\) y \(\theta=\operatorname{atan2}(b,a)\). Es
exactamente la identificación \(\rho e^{i\theta}\leftrightarrow\rho R(\theta)\) de la
clase 1: la matriz de la izquierda tiene dos parámetros, no cuatro. Leer \(a=x_u\),
\(b=y_u\) y comparar con la jacobiana general
\(\begin{pmatrix}x_u&x_v\\y_u&y_v\end{pmatrix}\) da Cauchy--Riemann.

(i)\(\Leftrightarrow\)(ii), por la SVD. Si \(J=U\Sigma V^T\) con \(\Sigma=kI\), entonces
\(J=k\,UV^T\), y \(UV^T\) es ortogonal (producto de ortogonales) con
\(\det(UV^T)=\det J/k^2>0\), es decir, una rotación: \(J=k\,R(\theta)\). Al revés, si
\(J=\rho R(\theta)\), su SVD es \(U=R(\theta)\), \(\Sigma=\rho I\), \(V=I\), y los dos
valores singulares valen \(\rho\). Geométricamente: la circunferencia unidad va a una
elipse de semiejes \(\sigma_1,\sigma_2\); la elipse es circunferencia si y sólo si los
semiejes son iguales.

(i)\(\Leftrightarrow\)(iii), por el producto \(J^TJ\). La primera forma inducida es
\(ds^2_{\text{mapa}}=(du\ dv)\,J^TJ\,(du\ dv)^T\), de modo que \(E,F,G\) son las
entradas de \(J^TJ\). Si \(J=\rho R\), \(J^TJ=\rho^2R^TR=\rho^2I\): \(F=0\) y
\(E=G=\rho^2\). Al revés, si \(J^TJ=k^2I\), la matriz \(J/k\) cumple \((J/k)^T(J/k)=I\),
es ortogonal; si además \(\det J>0\) es una rotación y \(J=k\,R(\theta)\). (Con
\(\det J<0\) sería una reflexión por una homotecia: conserva los ángulos pero invierte la
orientación, una similaridad inversa; las proyecciones cartográficas se toman directas.)
\(\square\)

\paragraph{Por qué (iii) es ``conservar ángulos''.} El ángulo entre dos pasitos \(p,q\)
del terreno se calcula con el producto interno, \(\cos\alpha=p\cdot q/(\|p\|\|q\|)\). En el
mapa los pasitos son \(Jp\) y \(Jq\), y \((Jp)\cdot(Jq)=p^T(J^TJ)q=k^2\,p\cdot q\): el
numerador se multiplica por \(k^2\), cada norma por \(k\), y el coseno no cambia. Es el
mismo argumento de la clase 1 (el producto interno es invariante por rotaciones) con una
homotecia más. Y es lo único que lo logra: si \(J^TJ\) no fuera múltiplo de la identidad,
sus dos autovectores (las direcciones \(v_1,v_2\) de la SVD) se estirarían distinto y
alguna dirección intermedia giraría respecto de las otras.

\paragraph{Cuánto se deforma un ángulo cuando no es conforme.} Con la SVD podemos
medirlo. En las bases \(v_1,v_2\) y \(u_1,u_2\) la jacobiana es \(\operatorname{diag}(\sigma_1,\sigma_2)\),
\(\sigma_1\ge\sigma_2>0\). La dirección de ángulo \(\alpha\) del terreno,
\((\cos\alpha,\sin\alpha)\), va a \((\sigma_1\cos\alpha,\sigma_2\sin\alpha)\), de ángulo
\(\alpha'\) con \(\tan\alpha'=r\tan\alpha\), \(r=\sigma_2/\sigma_1\le1\). El giro de esa
dirección es \(\delta=\alpha-\alpha'\), con
\[
  \tan\delta=\frac{\tan\alpha-\tan\alpha'}{1+\tan\alpha\tan\alpha'}
  =\frac{(1-r)\,t}{1+r\,t^2},\qquad t=\tan\alpha .
\]
Derivando en \(t\) y anulando, \(1+rt^2-2rt^2=0\), o sea \(t^2=1/r\): el giro es máximo
en la dirección con \(\tan\alpha=\sqrt{\sigma_1/\sigma_2}\), y ahí
\(\tan\alpha'=\sqrt{\sigma_2/\sigma_1}=1/\tan\alpha\), de modo que \(\alpha+\alpha'=90^\circ\):
la dirección y su imagen son simétricas respecto de la diagonal de la elipse. En ese
máximo \(\tan\delta=(1-r)/(2\sqrt r)\), y pasando a seno
(\(\sin\delta=\tan\delta/\sqrt{1+\tan^2\delta}\), con \((1-r)^2+4r=(1+r)^2\)):
\[
  \sin\delta_{\max}=\frac{1-r}{1+r}=\frac{\sigma_1-\sigma_2}{\sigma_1+\sigma_2}.
\]
La dirección simétrica, de ángulo \(-\alpha\), gira lo mismo en sentido contrario, así
que el ángulo \(2\alpha\) entre ambas es el que más cambia, y cambia en
\[
  \boxed{\ \omega=2\arcsin\frac{\sigma_1-\sigma_2}{\sigma_1+\sigma_2}\ }
\]
(deformación angular máxima de Tissot). Control: \(\omega=0\) si y sólo si
\(\sigma_1=\sigma_2\), que es (ii). Con \(r=1/2\), \(\omega=38{,}94^\circ\), lo mismo
que da maximizar numéricamente \(2(\alpha-\alpha')\).

\paragraph{Las proyecciones del cierre del apunte del elipsoide, con este lenguaje.} En
la sección~\ref{sec:arco-meridiano} planteamos las cinco proyecciones que el apunte del
profesor deja como ``puntapié inicial'' (página 20) y resolvimos las cilíndricas
integrando una condición cada vez. Ahora las vemos con las cuatro caras. Sobre el
elipsoide, con la base \((\varphi,\lambda)\),
\[
  ds^2=M^2\,d\varphi^2+N^2\cos^2\varphi\,d\lambda^2 :
\]
\(F=0\), la grilla de meridianos y paralelos es ortogonal, pero un paso \(d\varphi=d\lambda\)
no es un cuadradito: mide \(M\,d\varphi\) hacia el norte y \(N\cos\varphi\,d\lambda\) hacia
el este. En una proyección cilíndrica \(x=a\lambda\), \(y=y(\varphi)\), la jacobiana en las
bases ortonormales del terreno y del plano es diagonal,
\[
  J=\operatorname{diag}(h,k),\qquad
  h=\frac{dy/d\varphi}{M}\ \ (\text{meridiano}),\qquad
  k=\frac{dx/d\lambda}{N\cos\varphi}=\frac{a}{N\cos\varphi}\ \ (\text{paralelo}),
\]
así que \(\{\sigma_1,\sigma_2\}=\{h,k\}\) y cada condición es una ecuación en \(h\) y \(k\).
\begin{enumerate}[label=(\alph*)]
  \item \textbf{Mercator: conforme.} Cara (ii), \(h=k\):
  \[
    \frac{dy}{d\varphi}=\frac{a\,M}{N\cos\varphi}
    \quad\Longrightarrow\quad
    y=a\,q(\varphi),\qquad
    q(\varphi)=\int_0^\varphi\frac{M}{N\cos u}\,du ,
  \]
  la \emph{latitud isométrica}. La integral sale con una identidad: como
  \(M/N=(1-e^2)/(1-e^2\sin^2\varphi)\),
  \[
    \frac{M}{N\cos\varphi}=\frac{1-e^2}{(1-e^2\sin^2\varphi)\cos\varphi}
    =\frac1{\cos\varphi}-\frac{e^2\cos\varphi}{1-e^2\sin^2\varphi}
  \]
  (se verifica poniendo a común denominador: \(1-e^2\sin^2\varphi-e^2\cos^2\varphi=1-e^2\)).
  El primer término integra \(\ln\tan(\pi/4+\varphi/2)\), el de Mercator esférico. En el
  segundo, con \(t=\sin\varphi\), queda \(e^2\int dt/(1-e^2t^2)\), y las fracciones simples
  \(\frac1{1-e^2t^2}=\frac12\bigl(\frac1{1+et}+\frac1{1-et}\bigr)\) dan
  \(\frac e2\ln\frac{1+et}{1-et}\). Por lo tanto
  \[
    \boxed{\ q(\varphi)=\ln\tan\Bigl(\frac\pi4+\frac\varphi2\Bigr)
    -\frac e2\ln\frac{1+e\sin\varphi}{1-e\sin\varphi}\ } .
  \]
  Verificación derivando: \(\frac{d}{d\varphi}\ln\tan(\pi/4+\varphi/2)=1/\cos\varphi\) y
  \(\frac{d}{d\varphi}\frac e2\ln\frac{1+e\sin\varphi}{1-e\sin\varphi}
  =\frac e2\Bigl(\frac{e\cos\varphi}{1+e\sin\varphi}+\frac{e\cos\varphi}{1-e\sin\varphi}\Bigr)
  =\frac{e^2\cos\varphi}{1-e^2\sin^2\varphi}\), que es lo que había que obtener; y
  numéricamente, en \(\varphi=-34{,}6^\circ\), \(q=-0{,}640531\) tanto por la fórmula como
  por \texttt{quad}, y \(a\,q=-4\,085\,391{,}659\) m, el valor de la sección~\ref{sec:arco-meridiano}.
  Lo importante no es la fórmula sino lo que hace: en las coordenadas \((\lambda,q)\), con
  \(d\varphi=(N\cos\varphi/M)\,dq\), la métrica del elipsoide queda
  \[
    ds^2=N^2\cos^2\varphi\,\bigl(dq^2+d\lambda^2\bigr),
  \]
  múltiplo escalar de la métrica plana: cara (iii) con \(k=N\cos\varphi\). La grilla
  \((\lambda,q)\) \emph{es} la carta conforme, y Mercator es simplemente dibujarla con
  \(x=a\lambda\), \(y=aq\).
  \item \textbf{Cilíndrica equivalente.} Conservar áreas es \(\det J=hk=1\):
  \[
    \frac{dy}{d\varphi}=\frac{MN\cos\varphi}{a}\quad\Longrightarrow\quad
    y=\frac1a\int_0^\varphi MN\cos u\,du ,
  \]
  que es la \(\tfrac a2\,q_A(\varphi)\) de la sección~\ref{sec:arco-meridiano} (ahí la
  función se llamó \(q\); acá la escribimos \(q_A\), de ``autálica'', para no confundirla
  con la latitud isométrica). Los ángulos no se conservan: \(h=1/k\neq k\) salvo en el
  ecuador.
  \item \textbf{Cilíndrica equidistante.} Meridianos a escala verdadera, \(h=1\):
  \(y=S(\varphi)\), el arco de meridiano de la sección~\ref{sec:arco-meridiano}. Ni
  conforme ni equivalente: \(\sigma_1=k\), \(\sigma_2=1\).
  \item \textbf{Esfera autálica.} El área del elipsoide se integra con la misma jacobiana
  \(3\times2\) de la sección~\ref{sec:metrica}: el elemento de área es
  \(\sqrt{EG-F^2}\,d\varphi\,d\lambda=MN\cos\varphi\,d\varphi\,d\lambda\), y
  \[
    A=\int_{-\pi/2}^{\pi/2}2\pi N\cos\varphi\,M\,d\varphi
    =2\pi a^2(1-e^2)\int_{-1}^{1}\frac{dt}{(1-e^2t^2)^2}
  \]
  \[
    \phantom{A}=2\pi a^2(1-e^2)\Bigl[\frac{t}{2(1-e^2t^2)}+\frac1{4e}\ln\frac{1+et}{1-et}\Bigr]_{-1}^{1},
  \]
  con \(t=\sin\varphi\) y \(MN\cos\varphi=a^2(1-e^2)\cos\varphi/(1-e^2\sin^2\varphi)^2\)
  (la primitiva se verifica derivando, igual que en (a)). Evaluando,
  \[
    \boxed{\ A=2\pi a^2+\pi\frac{b^2}{e}\ln\frac{1+e}{1-e}
    =2\pi a^2\Bigl[1+\frac{1-e^2}{e}\operatorname{artanh}e\Bigr]\ },
    \qquad R_A=\sqrt{\frac{A}{4\pi}} .
  \]
  En GRS80, \(A=510\,065\,621{,}7\) km\(^2\) y \(R_A=6\,371\,007{,}181\) m (la integral
  numérica da lo mismo con nueve cifras). Es el \(R_q\) de la sección~\ref{sec:arco-meridiano};
  para comparar, el radio medio \((2a+b)/3=6\,371\,008{,}771\) m y el de igual volumen
  \(\sqrt[3]{a^2b}=6\,371\,000{,}790\) m: los tres ``radios de la Tierra'' difieren en
  metros. Pasar a la esfera autálica es equivalente (cara \(\det J=1\)), no conforme.
  \item \textbf{Transversa de Mercator (Gauss--Krüger), el planteo.} Queremos conforme
  \emph{y} el meridiano central \(\lambda_0\) sin deformación (\(h=k=1\) sobre él), porque
  la Argentina se cubre con fajas de \(3^\circ\) de ancho (POSGAR) donde lo que importa es
  deformar poco cerca del meridiano central, no cerca del ecuador. Con las cuatro caras la
  estrategia es clara: la grilla \((q,\lambda)\) ya es conforme por (a), y por (iv) toda
  función analítica de \(z=q+i(\lambda-\lambda_0)\) también lo es. Buscamos entonces
  \[
    w=\mathrm N+i\,\mathrm E=f(z),\qquad f \text{ analítica},\qquad
    f(q)=S(\varphi(q))\ \text{para } \lambda=\lambda_0 ,
  \]
  con \(S\) el arco de meridiano: sobre el meridiano central el norte es la distancia
  medida y el este es cero. Esa condición fija \(f\) sobre el eje real, y una función
  analítica queda determinada por sus valores en una recta (es la unicidad de la
  prolongación analítica): la carta es única. Para escribirla, desarrollamos en serie de
  Taylor en \(\Delta\lambda=\lambda-\lambda_0\) alrededor del eje real, con
  \('=d/dq\):
  \[
    f(q+i\Delta\lambda)=S+i\,\Delta\lambda\,S'-\frac{\Delta\lambda^2}2S''
    -i\frac{\Delta\lambda^3}6S'''+\cdots
  \]
  Las derivadas respecto de \(q\) son derivadas respecto de \(\varphi\) por
  \(d\varphi/dq=N\cos\varphi/M\): \(S'=M\cdot N\cos\varphi/M=N\cos\varphi\), y
  \(S''=\frac{N\cos\varphi}{M}\frac{d(N\cos\varphi)}{d\varphi}=-N\cos\varphi\sin\varphi\),
  porque \(\frac{d}{d\varphi}(N\cos\varphi)=-M\sin\varphi\) (sale derivando
  \(a\cos\varphi(1-e^2\sin^2\varphi)^{-1/2}\); lo controlamos numéricamente en
  \(-34{,}6^\circ\): \(3\,609\,219{,}47\) por los dos lados). Separando parte real e
  imaginaria, los dos primeros términos de cada coordenada son
  \[
    \mathrm N\approx S(\varphi)+\frac{\Delta\lambda^2}2\,N\sin\varphi\cos\varphi,\qquad
    \mathrm E\approx N\cos\varphi\,\Delta\lambda
    +\frac{\Delta\lambda^3}{6}\,N\cos^3\varphi\,(1-\tan^2\varphi+e'^2\cos^2\varphi),
  \]
  y la escala \(k=|f'(z)|/(N\cos\varphi)\approx1+\tfrac12\Delta\lambda^2\cos^2\varphi\)
  (es \(1\) sobre el meridiano central y crece con el cuadrado de la distancia a él; por
  eso Gauss--Krüger argentino no aplica factor de reducción y las fajas son angostas). El
  término cúbico de \(\mathrm E\) lo transcribimos de la bibliografía, no lo deducimos:
  sale de \(S'''\) con la misma mecánica. Control con los dos primeros términos, en
  \(\varphi=-34{,}6^\circ\) y \(\Delta\lambda=1{,}5^\circ\) (el borde de una faja):
  \(\mathrm N=-3\,831\,240{,}9\) m, \(\mathrm E=137\,595{,}3\) m, contra
  \(-3\,831\,241{,}1\) y \(137\,600{,}9\) de \texttt{pyproj} (\texttt{tmerc}, GRS80):
  \(0{,}2\) m en el norte y \(5{,}6\) m en el este, que es el término cúbico que falta.
  Queda como desarrollo complementario; lo que la materia aporta es el planteo:
  \emph{conforme} es ``analítica'', y ``sin deformación sobre el meridiano central'' es la
  condición de contorno que la vuelve única.
\end{enumerate}

\begin{resultado}[Buenos Aires y Ushuaia en las tres cilíndricas, GRS80]
\(\varphi=-34{,}6^\circ\) (Buenos Aires) y \(\varphi=-54^\circ\) (Ushuaia). Latitud
isométrica \(q\), escalas \(h\) (meridiano) y \(k=a/(N\cos\varphi)\) (paralelo),
deformación angular máxima \(\omega=2\arcsin\frac{|h-k|}{h+k}\) y factor de área \(hk\).

{\small
\begin{center}
\begin{tabular}{l l rrrr}
\toprule
\(\varphi\) & Proyección & \(h\) & \(k\) & \(\omega\) & \(hk\)\\
\midrule
\(-34{,}6^\circ\) & Mercator (\(h=k\)) & 1{,}2136 & 1{,}2136 & \(0^\circ\) & 1{,}4727\\
\(q=-0{,}6405\) & Equidistante (\(h=1\)) & 1{,}0000 & 1{,}2136 & \(11{,}0725^\circ\) & 1{,}2136\\
 & Equivalente (\(hk=1\)) & 0{,}8240 & 1{,}2136 & \(22{,}0423^\circ\) & 1{,}0000\\
\midrule
\(-54^\circ\) & Mercator (\(h=k\)) & 1{,}6976 & 1{,}6976 & \(0^\circ\) & 2{,}8817\\
\(q=-1{,}1188\) & Equidistante (\(h=1\)) & 1{,}0000 & 1{,}6976 & \(29{,}9731^\circ\) & 1{,}6976\\
 & Equivalente (\(hk=1\)) & 0{,}5891 & 1{,}6976 & \(57{,}9945^\circ\) & 1{,}0000\\
\bottomrule
\end{tabular}
\end{center}
}

Lectura: en Mercator un ángulo medido en Ushuaia se dibuja exacto, pero cada metro del
terreno son \(1{,}70\) en el mapa y cada hectárea \(2{,}88\); en la equivalente las áreas
son exactas pero hay direcciones que giran \(29^\circ\): un ángulo de \(119^\circ\) entre dos
alineaciones se dibuja de \(61^\circ\) (las dos direcciones más giradas y sus imágenes son
simétricas respecto de la diagonal, así que el ángulo y su imagen son suplementarios). No hay
proyección con las dos columnas en \(1\): es el \emph{theorema egregium} en números.
\end{resultado}

\paragraph{Control numérico.} Los valores de la tabla, \(q\), \(A\) y \(R_A\) los
calculamos con NumPy y SciPy (\texttt{quad}) en GRS80, \(a=6\,378\,137\) m,
\(f=1/298{,}257222101\), \(e^2=0{,}00669438\): \(M=a(1-e^2)/(1-e^2\sin^2\varphi)^{3/2}\),
\(N=a/(1-e^2\sin^2\varphi)^{1/2}\), \(k=a/(N\cos\varphi)\), y \(h=k\), \(1\), \(1/k\)
según la proyección. En \(-54^\circ\): \(M=6\,377\,306{,}911\) m, \(N=6\,392\,156{,}080\) m,
\(y_{\text{Mercator}}=aq=-7\,135\,562{,}567\) m, \(S=-5\,985\,916{,}028\) m,
\(y_{\text{equivalente}}=-5\,140\,508{,}988\) m. La \(q\) de la fórmula cerrada y la de
la integral numérica coinciden en las seis cifras mostradas, y
\(dq/d\varphi=1{,}697349\) por diferencias finitas, igual a \(M/(N\cos\varphi)\).
